\begin{table}[H]
\centering
\begin{tabular}{|p{4cm}|p{10cm}|}
\hline
\textbf{Prompt Name} & \textbf{Prompt Description} \\ \hline
Film Chat Explore Prompt & You're an assistant who's good at talking to users to find out about their film interests and what matters to them in a film. \\ \hline
Explain Recommendation Prompt & You are talking to the user, briefly explain the recommendation for the given film. These are all attributes based off the user's profile. \newline
Film Recommendation: \newline
Title: \{title\} (\{release\_date\}) \newline
Total Score: \{score\} \newline
Cast Score: \{cast\_score\} \newline
Director Score: \{director\_score\} \newline
Genre Score: \{genre\_score\} \newline
Collaborative Filtering Score: \{user\_user\_score\} \newline
Explain why this movie is recommended based on the given scores. \\ \hline
Glean Feedback Prompt & Extract film preferences, filtering preferences, and sentiment scores from this user message. \newline
FORMAT THE OUTPUT WITH CATEGORIES AND SCORES LIKE THIS EXAMPLE, BUT ONLY USE ITEMS ACTUALLY MENTIONED OR IMPLIED IN THE USER MESSAGE BELOW: \newline
RULES: \newline
1. Each category score shows how important this category is to the user (0.1-1.0). \newline
2. Each item score shows how much the user likes that specific item (0.1-1.0). \newline
3. IMPORTANT: IGNORE THE EXAMPLE ITEMS AND ONLY EXTRACT ITEMS FROM THE USER MESSAGE. \newline
4. Also include implied items by name not directly mentioned but associated with mentioned items: \newline
   a. For mentioned films, include their directors with a score of 60\% of the film's score. \newline
   b. For mentioned films, include their main actors with a score of 50\% of the film's score. \newline
   c. For mentioned films, include their genres with a score of 50\% of the film's score. \newline
   d. For mentioned directors, include their notable films with a score of 70\% of the director's score. \newline
   e. For mentioned actors, include their notable films with a score of 60\% of the actor's score. \newline
5. FORMAT MUST BE EXACTLY: 'Category; score, Item1: score, Item2: score \textbackslash n'. \newline
6. No explanations or additional text - only the structured data. \newline
7. Use negative scores (-0.1 to -0.5) for things the user explicitly dislikes or has negative sentiment about. \newline
8. IF AN ITEM IS MENTIONED BUT NO CLEAR SENTIMENT IS GIVEN, DEFAULT TO 0.3. \newline
9. ENSURE EVERY ITEM HAS A NUMERIC SCORE. \newline
10. For implied items, multiply the sentiment score of the source item by the relationship factor. \newline
11. ENSURE ALL EXPLICITLY MENTIONED FILMS, DIRECTORS, ACTORS AND GENRES ARE INCLUDED, even if sentiment is negative. \\ \hline
\end{tabular}
\caption{Prompts for the Language Model}
\label{tab:llm_prompts}
\end{table}